\paragraph{Processing pipeline overview.}
\label{sec:pipeline_overview}

The proposed method is an epoch-aware pipeline that converts a continuous multichannel EEG recording into channel-wise blink-region detections through five sequential operations. The first two operations prepare the signal, in which (1) the EEG is downsampled to 100 Hz and band-pass filtered between 1 and 20 Hz, and (2) segmented into fixed-length, non-overlapping 30 s epochs. The remaining three operations constitute the detection method. (3) Stage A applies an epoch-level, multichannel peak-to-peak screening with a union rule to identify candidate epochs containing prominent amplitude excursions. (4) Stage B pools the samples from these candidate epochs to estimate a robust sample-level threshold independently for each channel. (5) Stage C applies each channel's threshold to the corresponding channel across all valid epochs to localize blink regions, each described by its epoch, channel, onset time, and duration. Figure~\ref{fig:pipeline_flow} summarizes the complete processing pipeline, while Stages A, B, and C are described in detail in the following subsections.


\begin{figure*}[!t]
    \centering
    \includegraphics[width=\textwidth]{flowchart}
    \caption{
        Overview of the proposed epoch-aware pipeline for EEG blink-region detection. Continuous multichannel EEG signals are downsampled to 100 Hz, band-pass filtered from 1 to 20 Hz, and segmented into non-overlapping 30-s epochs. Stage A applies multichannel peak-to-peak screening with a union rule to select candidate epochs, represented by E2, E5, and E7 in this example. Stage B pools the samples from these candidate epochs to estimate a robust sample-level threshold independently for each channel
        , illustrated here for $FP1$.
        Stage C applies each channel's threshold to the corresponding channel across all valid epochs, including non-candidate epochs, and reports each detected blink region by its epoch, channel, onset time, and duration. Pale yellow shading denotes Stage A candidate epochs, red shading denotes detected blink regions, and dashed green lines denote the channel-specific thresholds.
    }
    \label{fig:pipeline_flow}
\end{figure*}

\paragraph{Preprocessing}
First, the continuous recording is downsampled to a common sampling grid of \(100\)~Hz. Standardising the sampling rate makes the
amplitude thresholds and the minimum-duration gate comparable across recordings acquired on different hardware, and reduces the per-sample computational cost of the downstream stages without discarding any information needed to resolve a blink, whose
spectral content lies well below the new Nyquist frequency.
%
Second, a zero-phase band-pass filter from \(1.0\)~Hz  to \(20.0\)~Hz
is applied which is removes slow drifts and DC offsets at the low end and high-frequency muscle and line noise at the high end, isolating the low-frequency, high-amplitude deflection that characterises an eye blink.
%
Third, each filtered recording is split into non-overlapping, fixed-length epochs. Epoching is what makes the thresholds
adaptive, rather than deriving a single level for the entire session, the pipeline treats each short epoch as a local statistical context.

